\section*{Capítulo \ref{ch:b3:lp} --- Los espacios $L^p$}

\begin{solution}{exo:b3:lp:1}
(a) $p = q = 2$ en el~\cref{thm:b3:lp:holder}: $\abs{\int f\bar
g\,\dd\mu} \leq \int\abs{fg} \leq \norm f_2\norm g_2$ ---
Cauchy--Schwarz, con igualdad si y solo si $\abs f, \abs g$ son
proporcionales y las fases están alineadas.

(b) Sobre un espacio de probabilidad, para $p \leq q$: aplíquese Jensen
(el~\cref{exo:b3:lp:10}) con la función convexa $\Phi(t) =
\abs t^{q/p}$ a la función $\abs f^p$:
$\bigl(\int\abs f^p\bigr)^{q/p} \leq \int\abs f^q$, i.e.\
$\norm f_p \leq \norm f_q$.

(c) $\int\abs{fg} = \norm f_1\norm g_\infty$ si y solo si $\abs{g} =
\norm g_\infty$ en casi todo punto de $\{f \neq 0\}$ (la desigualdad
$\abs{fg} \leq \abs f\norm g_\infty$ ha de ser una igualdad en casi todo
punto).
\end{solution}

\begin{solution}{exo:b3:lp:2}
$\int_0^1 x^{-p/2}\dd x < \infty$ si y solo si $p < 2$: la primera está
en $L^p$ para $p \in \intco12$. $\int_1^\infty x^{-p/2}\dd x <
\infty$ si y solo si $p > 2$: la segunda, para $p \in \intoo2\infty$ (y
$p = \infty$: está acotada --- inclúyase). Tercera: para $p <
2$, dominada por $x^{-p/2}$: \omterm{def:b3:lebesgue:l1}{integrable};
para $p = 2$, sustitúyase $u = -\ln x$:
$\int_0^1\frac{\dd x}{x(1 + \abs{\ln
x})^2} = \int_0^\infty\frac{\dd u}{(1 + u)^2} < \infty$; para $p > 2$ la
potencia domina: divergente. Luego $p \in \intcc12$. Cuarta:
$\int_\R\frac{\dd x}{(1 + \abs x)^p} < \infty$ si y solo si $p >
1$; y está acotada, luego también $L^\infty$: $p \in
\intoc1\infty$. Moraleja: la \omterm{def:b3:lebesgue:l1}{integrabilidad}
en $0$ prefiere $p$ pequeño y en $\infty$ prefiere $p$ grande;
combinando ambas obstrucciones, no hay inclusión entre espacios
$L^p(\R)$.
\end{solution}

\begin{solution}{exo:b3:lp:3}
(a) $\norm{f_n}_p^p = \lambda(I_n) \to 0$ (en el nivel diádico $k$ la
longitud es $2^{-k}$). Pero todo $x$ está en un intervalo de cada nivel
diádico: $f_n(x) = 1$ infinitas veces y $= 0$ infinitas veces
(intervalos del mismo nivel que no contienen $x$): no hay convergencia
en ningún punto.

(b) $f_{n_k} = \mathbf 1_{\intcc0{2^{-k}}}$ (el primer intervalo de cada
nivel) converge a $0$ en todo $x > 0$: en casi todo punto.

(c) En casi todo punto pero no en $L^1$:
$n\mathbf 1_{\intoo0{1/n}} \to 0$ en casi todo punto, con integral $1$.
En $L^1$ pero en ningún $L^p$ ($p > 1$): $g_n =
\eu^n\,\mathbf 1_{(0,\ \eu^{-n}/n)}$: $\norm{g_n}_1 = \frac1n
\to 0$, mientras que $\norm{g_n}_p^p = \eu^{(p-1)n}/n \to \infty$ para
todo $p > 1$.
\end{solution}

\begin{solution}{exo:b3:lp:4}
Superior: $\norm f_p \leq \mu(X)^{1/p}\norm f_\infty$
(la~\cref{prop:b3:lp:inclusions}(a) con $q = \infty$), y
$\mu(X)^{1/p} \to 1$. Inferior: para $\varepsilon > 0$, $A =
\{\abs f > \norm f_\infty - \varepsilon\}$ cumple $\mu(A) > 0$
(definición del supremo esencial), y
\[
\norm f_p \geq \Bigl(\int_A \abs f^p\Bigr)^{1/p}
\geq (\norm f_\infty - \varepsilon)\,\mu(A)^{1/p}
\xrightarrow[p\to\infty]{} \norm f_\infty - \varepsilon .
\]
\end{solution}

\begin{solution}{exo:b3:lp:5}
(a) Dos veces: la conversión $\norm{\tau_hg - g}_p \leq
C^{1/p}\norm{\tau_hg - g}_\infty$ (soporte de
\omterm{def:b3:measure:measure}{medida} finita) degenera para
$p = \infty$ únicamente en que allí falla la
\emph{\omterm{ex:b3:lebesgue:gamma}{densidad} de $\mathcal C_c$} --- esa
es la laguna real: el paso (2) del~\cref{thm:b3:lp:density} no tiene
análogo $L^\infty$.

(b) Si $f$ tiene un representante uniformemente
\omterm{def:b3:topology:continuity}{continuo} $g$:
$\norm{\tau_hf - f}_\infty = \sup_x\abs{g(x - h) - g(x)} \to
0$. Recíprocamente, supóngase $\norm{\tau_hf - f}_\infty \to 0$. Las
regularizaciones $f_\varepsilon = f * \rho_\varepsilon$ son
\omterm{def:b3:topology:continuity}{continuas}, y
\[
\norm{f_\varepsilon - f}_\infty
\leq \sup_{\norm y \leq \varepsilon}\norm{\tau_yf -
f}_\infty \longrightarrow 0
\]
(la convolución es un promedio de trasladadas). Cada $f_\varepsilon$ es
uniformemente \omterm{def:b3:topology:continuity}{continua}
($\norm{\tau_hf_\varepsilon - f_\varepsilon}_\infty \leq
\norm{\tau_hf - f}_\infty$, promediando), y un límite uniforme de
funciones uniformemente \omterm{def:b3:topology:continuity}{continuas}
lo es: $f$ coincide en casi todo punto con una función uniformemente
\omterm{def:b3:topology:continuity}{continua}.
\end{solution}

\begin{solution}{exo:b3:lp:6}
Por Hölder, para todo $x$ el integrando $y \mapsto f(x -
y)g(y)$ está en $L^1$ con $\abs{(f*g)(x)} \leq \norm f_p\norm
g_q$: definida en todo punto y acotada.
\omterm{def:b3:topology:continuity}{Continuidad} uniforme ($p <
\infty$):
\[
\abs{(f*g)(x + h) - (f*g)(x)}
= \Bigl|\int\bigl(\tau_{-h}f - f\bigr)(x - y)\,g(y)\dd y\Bigr|
\leq \norm{\tau_{-h}f - f}_p\,\norm g_q
\xrightarrow[h\to0]{} 0,
\]
uniformemente en $x$ (el~\cref{thm:b3:lp:density}(3)). Si $p =
\infty$, entonces $q = 1$: escríbase $f * g = g * f$ y repítase la misma
cota con la traslación actuando sobre $g \in L^1$.
\end{solution}

\begin{solution}{exo:b3:lp:7}
Fíjese $\chi \in \mathcal C_c^\infty(\intoo ab)$ con $\int\chi =
1$ y póngase $c = \int f\chi$. Sea $\varphi \in \mathcal
C_c^\infty$ arbitraria y $\psi = \varphi -
\bigl(\int\varphi\bigr)\chi$: entonces $\int\psi = 0$, de modo que
$\Phi(x)
= \int_a^x\psi$ define $\Phi \in \mathcal C_c^\infty(\intoo
ab)$ (se anula cerca de ambos extremos: cerca de $a$ trivialmente, y
cerca de $b$ porque la integral total es $0$) con $\Phi' = \psi$. La
hipótesis da $\int f\psi = \int f\Phi' = 0$, luego
\[
\int f\varphi = \Bigl(\int\varphi\Bigr)\int f\chi = \int
c\,\varphi
\quad\text{para todo } \varphi:
\qquad \int(f - c)\varphi = 0 .
\]
Por el~\cref{cor:b3:lp:fundlemma} (localizado en $\intoo ab$), $f =
c$ en casi todo punto.
\end{solution}

\begin{solution}{exo:b3:lp:8}
Sean $3\delta < d(K, \R^d\setminus U)$ (positiva:
el~\cref{exo:b3:topology:6}(b); si $U = \R^d$, sirve cualquier
$\delta$), $K_\delta = \{x : d(x, K) \leq \delta\}$ y
\[
\varphi = \mathbf 1_{K_\delta} * \rho_{\delta/2} .
\]
Entonces $\varphi \in \mathcal C^\infty$
(el~\cref{thm:b3:lp:regularization}(1); el indicador es $L^1$),
$0 \leq \varphi \leq 1$ ($\int\rho = 1$), $\varphi = 1$ sobre $K$ (para
$x \in K$, $\bar B(x, \delta/2) \subseteq K_\delta$, de modo que la
convolución integra $\rho$ por completo) y
$\operatorname{supp}\varphi \subseteq K_{3\delta/2} \subseteq
U$: de soporte \omterm{def:b3:topology:compact}{compacto} ($K_\delta$ está acotado). Partición de la
unidad: dado $K \subseteq U_1\cup\dots\cup U_m$, elíjanse (por
\omterm{def:b3:topology:compact}{compacidad}) \omterm{def:b3:topology:compact}{compactos}
$K_i \subseteq U_i$ con $K \subseteq
\bigcup \mathring K_i$, tómense $\varphi_i$ como arriba para $(K_i,
U_i)$ y póngase $\psi_i = \varphi_i\prod_{j <i}(1 -
\varphi_j)$: cada $\psi_i \in \mathcal C_c^\infty(U_i)$, y
$\sum_i\psi_i = 1 - \prod_i(1 - \varphi_i) = 1$ sobre $K$.
\end{solution}

\begin{solution}{exo:b3:lp:9}
(a) El exponente de interpolación de $(p_0, q_0) = (1, \infty)$ en
$r = p$ es $\alpha = \frac1p$: la~\cref{prop:b3:lp:inclusions}(c) da
$\norm f_p \leq \norm
f_1^{1/p}\norm f_\infty^{1 - 1/p}$.

(b) Tomando logaritmos en la~\cref{prop:b3:lp:inclusions}(c):
$\ln\norm f_r \leq
\alpha\ln\norm f_p + (1-\alpha)\ln\norm f_q$, donde $\frac1r$ es la
misma combinación convexa de $\frac1p, \frac1q$: $\frac1p
\mapsto \ln\norm f_p$ es convexa. Ejemplo con pertenencia a $L^p$
exactamente en $\intoo{p_0}{p_1}$:
\[
f(x) = x^{-1/p_1}\,\mathbf 1_{\intoo01}(x) +
x^{-1/p_0}\,\mathbf 1_{\intco1\infty}(x):
\]
el primer término está en $L^p$ si y solo si $p < p_1$, y el segundo si
y solo si $p >
p_0$.
\end{solution}

\begin{solution}{exo:b3:lp:10}
Sea $m = \int f\,\dd\mu \in \R$. La convexidad proporciona una recta de
apoyo en $m$: existe $s$ con $\Phi(t) \geq \Phi(m) + s(t -
m)$ para todo $t$ (tómese $s$ entre las derivadas laterales, que existen
para las funciones convexas). Sustitúyase $t = f(x)$ e intégrese contra
la probabilidad $\mu$:
\[
\int\Phi\circ f\,\dd\mu \geq \Phi(m) + s\Bigl(\int f - m\Bigr)
= \Phi\Bigl(\int f\,\dd\mu\Bigr)
\]
(\omterm{def:b3:lebesgue:measurable}{medibilidad}: $\Phi$ es
\omterm{def:b3:topology:continuity}{continua}; y la
\omterm{def:b3:lebesgue:l1}{integrabilidad} de la parte negativa de
$\Phi\circ f$ la garantiza la recta de apoyo). Medias aritmética y
geométrica: sobre un conjunto finito con pesos $w_i$, tómense
$\Phi = \exp$ y $f = \sum(\ln a_i)\mathbf 1_i$:
$\exp\bigl(\sum w_i\ln a_i\bigr) \leq \sum w_ia_i$, i.e.\
$\prod a_i^{w_i} \leq \sum w_ia_i$. La monotonía de las normas es
\cref{exo:b3:lp:1}(b).
\end{solution}

\begin{solution}{exo:b3:lp:11}
(a) Supóngase primero $p, q, r < \infty$ y $f, g \geq 0$ (sustitúyanse
por sus valores absolutos). Los tres exponentes $r$,
$\alpha = \frac{pr}{r-p}$, $\beta = \frac{qr}{r-q}$
cumplen $\frac1r + \frac1\alpha + \frac1\beta = \frac1r +
\frac1p - \frac1r + \frac1q - \frac1r = 1$ (la relación de escala).
Descompóngase, para $x$ fijo,
\[
f(y)g(x-y) = \bigl[f(y)^pg(x-y)^q\bigr]^{1/r}\cdot
f(y)^{1 - p/r}\cdot g(x-y)^{1 - q/r},
\]
y Hölder con los tres exponentes da
\[
(f*g)(x) \leq \Bigl(\int f^pg(x-\cdot)^q\Bigr)^{1/r}
\norm f_p^{\,p/\alpha\cdot\alpha/p}\cdots
\]
con más precisión: el segundo factor es
$\bigl(\int f^{(1-p/r)\alpha}\bigr)^{1/\alpha} =
\norm f_p^{p(1/p - 1/r)}$ porque $(1 - \frac pr)\alpha = p$, y
análogamente el tercero es $\norm g_q^{q(1/q - 1/r)}$. Elévese a la
potencia $r$ e intégrese en $x$ (Tonelli sobre el primer factor):
\[
\norm{f*g}_r^r \leq \norm f_p^p\,\norm g_q^q\cdot
\norm f_p^{\,rp(1/p - 1/r)}\,\norm g_q^{\,rq(1/q - 1/r)}
= \norm f_p^r\,\norm g_q^r .
\]
Los casos límite ($r = \infty$, o algún exponente igual a su cota) son
Hölder simple o estimaciones directas.

(b) $r = \infty$ obliga a $q = p'$: $\abs{f*g(x)} \leq \norm
f_p\norm g_{p'}$ --- Hölder tras trasladar y reflejar. $q = 1$ da
$r = p$: $\norm{f*g}_p \leq \norm
g_1\norm f_p$, el caballo de batalla de la regularización (el motor
del~\cref{thm:b3:lp:regularization}). $p = q = 1$ da $r = 1$: el álgebra
de convolución
(\cref{thm:b3:product:convolution}).

(c) Sustitúyanse $f, g$ por $f_\lambda = f(\lambda\cdot)$,
$g_\lambda = g(\lambda\cdot)$: entonces $f_\lambda*g_\lambda =
\lambda^{-d}(f*g)(\lambda\cdot)$ y, comparando normas,
\[
\text{MI} \sim \lambda^{-d - d/r}, \qquad
\text{MD} \sim \lambda^{-d/p - d/q} :
\]
una desigualdad válida para todo $f, g$ obliga a que los dos exponentes
de escala coincidan, es decir, a $1 + \frac1r = \frac1p + \frac1q$
exactamente. Cualquier otra combinación muere en $\lambda \to 0$ o en
$\infty$: la relación de Young no es una comodidad, sino una ley de
escala.
\end{solution}

\begin{solution}{exo:b3:lp:12}
(a) Normalícese $\norm f_p = \norm g_q = 1$. La demostración de Hölder
integra la desigualdad de Young $\abs{fg} \leq
\frac{\abs f^p}p + \frac{\abs g^q}q$; la igualdad de las integrales
obliga a la igualdad en casi todo punto en Young, lo cual (convexidad
estricta de $\exp$; igualdad si y solo si $a^p = b^q$) significa $\abs
f^p = \abs g^q$ en casi todo punto. Deshaciendo la normalización:
$\abs f^p\norm g_q^q = \abs g^q\norm f_p^p$ a.e.\ ---
proporcionalidad.

(b) Minkowski son dos Hölder aplicados a $\abs{f + g}^{p-1}
\abs f$ y $\abs{f+g}^{p-1}\abs g$; la igualdad obliga a la
proporcionalidad de (a) en ambos: $\abs f^p$ y $\abs g^p$ son cada una
proporcional a $\abs{f+g}^{(p-1)q} = \abs{f+g}^p$, de modo que
$\abs g = t\abs f$ en casi todo punto para cierta constante $t \geq 0$;
y la desigualdad triangular puntual inicial $\abs{f + g} \leq
\abs f + \abs g$ ha de ser también una igualdad en casi todo punto, lo
que para valores complejos significa que $f$ y $g$ tienen en casi todo
punto el mismo argumento allí donde ambas son no nulas. Combinando:
$g = tf$ en casi todo punto, con $t > 0$ (ambas no nulas).

(c) $p = 1$: hay igualdad en $\int\abs{f + g} = \int\abs f +
\int\abs g$ siempre que $f, g$ tengan el mismo patrón de signos (el
mismo argumento en casi todo punto) --- sin necesidad de
proporcionalidad: $f = \mathbf 1_{\intcc01}$ y
$g = \mathbf 1_{\intcc02}$ sirven. $p = \infty$:
$\norm{f+g}_\infty = \norm f_\infty +
\norm g_\infty$ en cuanto las dos funciones alcancen picos compatibles
en un punto común (o a lo largo de una sucesión común): basta $f = g\,$
cerca de un punto, con independencia del comportamiento en el resto. La
convexidad estricta de las bolas $L^p$ para $1 <
p < \infty$ --- y su fallo en los extremos --- es exactamente lo que
atestiguan estos casos de igualdad.
\end{solution}

\begin{solution}{pb:b3:lp:1}
\textbf{1.} $f$ tiene soporte en cierto $[\alpha, \beta]
\subseteq \intoo0\infty$, de modo que $F = 0$ sobre $[0, \alpha]$ y $F
\equiv F(\beta)$ sobre $[\beta, \infty)$: $Hf$ se anula cerca de $0$ y
es $O(1/x)$ en el infinito; $\int_1^\infty x^{-p}\dd x <
\infty$ para $p > 1$, y $Hf$ es
\omterm{def:b3:topology:continuity}{continua}: $Hf \in L^p$.

\textbf{2.} $\bigl(x^{1-p}F(x)^p\bigr)' = (1-p)x^{-p}F^p +
p\,x^{1-p}F^{p-1}f$. Ambos valores de frontera se anulan: en $0$ porque
$F = 0$ cerca de $0$; y en $\infty$ porque $x^{1-p}F^p \leq
F(\beta)^p x^{1-p} \to 0$ ($p > 1$). Integrando la identidad sobre
$\intoo0\infty$:
\[
0 = (1 - p)\int_0^\infty\Bigl(\frac Fx\Bigr)^p\dd x
+ p\int_0^\infty\Bigl(\frac Fx\Bigr)^{p-1}f(x)\,\dd x,
\]
que es la relación indicada.

\textbf{3.} Hölder con los exponentes $q = \frac p{p-1}$ y $p$:
\[
\int\Bigl(\frac Fx\Bigr)^{p-1}f
\leq \Bigl(\int\Bigl(\frac
Fx\Bigr)^{p}\Bigr)^{1 - 1/p}\,\norm f_p ,
\]
luego $\norm{Hf}_p^p \leq \frac p{p-1}\norm{Hf}_p^{p-1}\norm
f_p$; divídase (es finito por la pregunta 1, y si $0$ no hay nada que
demostrar).

\textbf{4.} Sean $f \in L^p$, $f \geq 0$. Elíjanse $\varphi_k \in
\mathcal C_c(\intoo0\infty)$ con $\varphi_k \to f$ en $L^p$
(el~\cref{thm:b3:lp:density}(2), intersecado con la semirrecta abierta
--- aproxímese $f\mathbf 1_{[1/k, k]}$ y diagonalícese) y sustitúyase
$\varphi_k$ por $\abs{\varphi_k}$ (que sigue siendo
\omterm{def:b3:topology:continuity}{continua} y está más cerca de
$f \geq 0$). Para todo $x > 0$ fijo, Hölder sobre $\intoo0x$ da
\[
\abs{H\varphi_k(x) - Hf(x)} \leq
\frac1x\,x^{1 - 1/p}\,\norm{\varphi_k - f}_p \to 0 :
\]
$H\varphi_k \to Hf$ puntualmente. Fatou y la pregunta 3:
\[
\int (Hf)^p \leq \liminf_k\int(H\varphi_k)^p
\leq \Bigl(\frac p{p-1}\Bigr)^p\liminf_k\norm{\varphi_k}_p^p
= \Bigl(\frac p{p-1}\Bigr)^p\norm f_p^p .
\]
Para $f$ con signo o compleja: $\abs{Hf} \leq H\abs f$ puntualmente, y
el caso no negativo se aplica a $\abs f$.

\textbf{5.} $\norm{f_A}_p^p = \int_1^A t^{-1}\dd t = \ln A$.
Para $1 \leq x \leq A$:
\[
(Hf_A)(x) = \frac1x\int_1^x t^{-1/p}\dd t
= \frac{x^{1 - 1/p} - 1}{(1 - \tfrac1p)\,x}
= \frac p{p-1}\,x^{-1/p}\bigl(1 - x^{-(1 - 1/p)}\bigr).
\]

\textbf{6.} Fíjense $\varepsilon > 0$ y $X_0$ con $(1 -
x^{-(1-1/p)})^p \geq 1 - \varepsilon$ para $x \geq X_0$. Entonces
\[
\norm{Hf_A}_p^p \geq \int_{X_0}^A\Bigl(\frac
p{p-1}\Bigr)^p\frac{1 - \varepsilon}{x}\,\dd x
= \Bigl(\frac p{p-1}\Bigr)^p(1 - \varepsilon)\,(\ln A - \ln
X_0),
\]
luego $\dfrac{\norm{Hf_A}_p^p}{\norm{f_A}_p^p} \geq \bigl(\frac
p{p-1}\bigr)^p(1 - \varepsilon)\bigl(1 - \frac{\ln X_0}{\ln
A}\bigr) \to \bigl(\frac p{p-1}\bigr)^p(1 - \varepsilon)$ cuando
$A \to \infty$: la constante no puede mejorarse.

\textbf{7.} La igualdad en la pregunta 3 obliga a la igualdad en Hölder:
$f^p$ proporcional a $\bigl(\frac Fx\bigr)^{(p-1)q}
= \bigl(\frac Fx\bigr)^p$ en casi todo punto, es decir,
$f = \gamma\,\frac Fx$ en casi todo punto para cierto $\gamma \geq 0$.
Como $F(x) = \int_0^xf$ es absolutamente
\omterm{def:b3:topology:continuity}{continua} con $F' = f$ en casi todo
punto, $F$ resuelve $F' =
\gamma F/x$: en todo intervalo donde $F > 0$, $(\ln F)' =
\gamma/x$, luego $F = c\,x^{\gamma}$ y $f = c\gamma
x^{\gamma - 1}$ allí. Pero ninguna potencia no nula $x^{\gamma-1}$
pertenece a $L^p(\intoo0\infty)$ (haría falta $p(\gamma - 1) < -1$ en
$\infty$ y $> -1$ en $0$: incompatible), y $F$ no puede anularse
idénticamente salvo que $f = 0$. Así que la igualdad exige $f =
0$.

\textbf{8.} Dada $(a_n)$ no creciente $\geq 0$, defínase la función
escalonada $g(t) = a_{\lceil t\rceil}$ sobre $\intoo0{+\infty}$: es no
creciente, con $\int_0^\infty g^p =
\sum_na_n^p$ y $\int_0^ng = a_1 + \dots + a_n$, luego
$(Hg)(n) = \frac{a_1 + \dots + a_n}n$. El promedio de una función no
creciente es no creciente, de modo que $Hg$ lo es, y
\[
\sum_{n\geq1}\Bigl(\frac{a_1{+}\dots{+}a_n}n\Bigr)^p
= \sum_{n\geq1}(Hg)(n)^p
\leq \sum_{n\geq1}\int_{n-1}^n (Hg)(t)^p\,\dd t
= \norm{Hg}_p^p
\leq \Bigl(\frac p{p-1}\Bigr)^p\sum_n a_n^p
\]
por la Parte I. Para una sucesión no negativa general, sea $(a_n^*)$ su
reordenación no creciente (posible cuando $a_n \to 0$, que podemos
suponer --- si no, ambos miembros son infinitos): el miembro derecho no
cambia, y cada suma parcial $a_1 + \dots
+ a_n$ es a lo sumo $a_1^* + \dots + a_n^*$ (los $n$ términos mayores):
el miembro izquierdo solo crece. De ahí la desigualdad para todo
$(a_n)$.

\textbf{9.} Si $(a_n) \in \ell^p$, la sucesión de medias de Cesàro está
en $\ell^p$ con norma $\leq \frac p{p-1}\norm
a_p$. Ejemplo ($p = 2$): $a_n = \frac1{\sqrt n\,\ln n}$ ($n
\geq 2$): $\sum a_n^2 = \sum\frac1{n\ln^2n} < \infty$, y sin embargo
$\sum a_n = \infty$ (criterio integral); Hardy sigue garantizando
$\sum_n\bigl(\frac{a_1 + \dots + a_n}n\bigr)^2 < \infty$.

\textbf{10.} Para $f = \mathbf 1_{\intcc01}$: $Hf(x) = 1$ en $\intoc01$
y $= \frac1x$ para $x \geq 1$: $\int_0^\infty Hf =
1 + \int_1^\infty\frac{\dd x}x = \infty$, mientras que $\norm f_1 =
1$. La demostración se hunde por dos sitios: la constante $\frac p{p-1}$
explota cuando $p \to 1$, y el término de frontera $x^{1-p}F^p$ ya no se
anula en el infinito para $p = 1$. La desigualdad de Hardy es un
fenómeno honestamente $p > 1$.

\textbf{11.} Tómese el intervalo más largo $B_{i_1}$; descártese todo
intervalo que lo corte; tómese el superviviente más largo $B_{i_2}$;
itérese (hay finitos intervalos). Los elegidos son disjuntos por
construcción, y todo $B$ descartado cortaba a un intervalo elegido al
menos igual de largo: un intervalo que corta a otro de longitud mayor o
igual está contenido en su triple, $B \subseteq
3B_{i_l}$.

\textbf{12.} $\{Mf > t\}$ es abierto: cada promedio $x \mapsto
\frac1{2r}\int_{x-r}^{x+r}\abs f$ es
\omterm{def:b3:topology:continuity}{continuo} (convergencia dominada en
$x$) y un supremo de funciones
\omterm{def:b3:topology:continuity}{continuas} es semicontinuo
inferiormente. Cada $x$ suyo posee $B_x =
\intoo{x-r_x}{x+r_x}$ con $\int_{B_x}\abs f >
t\,\lambda(B_x)$. Para $K \subseteq \{Mf > t\}$
\omterm{def:b3:topology:compact}{compacto}: un número finito de $B_x$
recubre $K$, Vitali (pregunta 11) extrae $B_1', \dots, B_k'$ disjuntos
con $K \subseteq
\bigcup_l3B_l'$, luego
\[
\lambda(K) \leq 3\sum_l\lambda(B_l') <
\frac3t\sum_l\int_{B_l'}\abs f \leq \frac3t\,\norm f_1
\]
por disjunción; la regularidad \omterm{def:b3:topology:interior}{interior}
(el~\cref{thm:b3:measure:regularity}) concluye.

\textbf{13.} Para $x > 1$: con $r \in \intcc{x-1}x$ el promedio es
$\frac{r - x + 1}{2r}$, creciente en $r$; para $r
\geq x$ es $\frac1{2r}$, decreciente: el supremo es $\frac1{2x}$,
alcanzado en $r = x$. Luego $M\mathbf 1_{\intcc01}
\notin L^1$. En general, si $\int_I\abs f = c > 0$ sobre un intervalo
acotado $I \subseteq \intcc{-C}C$, entonces $Mf(x) \geq
\frac{c}{2(\abs x + C)}$ para todo $x$: nunca
\omterm{def:b3:lebesgue:l1}{integrable} salvo que $f = 0$ en casi todo
punto.

\textbf{14.} Con $f_t = f\,\mathbf 1_{\abs f > t/2}$:
$M(f - f_t) \leq \frac t2$, luego $\{Mf > t\} \subseteq \{Mf_t
> \frac t2\}$ y la pregunta 12 da $\lambda(Mf > t) \leq
\frac6t\int_{\abs f > t/2}\abs f$. Fórmula de las capas
(la~\cref{prop:b3:product:layercake}) y Tonelli:
\[
\norm{Mf}_p^p = p\int_0^\infty t^{p-1}\lambda(Mf > t)\dd t
\leq 6p\int\abs f\int_0^{2\abs f}t^{p-2}\,\dd t\,\dd\lambda
= \frac{6p\,2^{p-1}}{p-1}\,\norm f_p^p .
\]

\textbf{15.} Escríbase $A_rf(x) = \frac1{2r}\int_{x-r}^{x+r}f$. Para $g$
\omterm{def:b3:topology:continuity}{continua}: $A_rg(x) \to g(x)$ en
todo punto. Dado $\varepsilon > 0$, descompóngase $f = g + h$ con $g$
\omterm{def:b3:topology:continuity}{continua} de soporte
\omterm{def:b3:topology:compact}{compacto} y $\norm h_1 < \varepsilon$
(el~\cref{thm:b3:lp:density}); entonces
\[
\limsup_{r\to0}\,\abs{A_rf(x) - f(x)} \leq Mh(x) +
\abs{h(x)},
\]
de modo que $\Omega_\delta = \{\limsup_r\abs{A_rf - f} > \delta\}
\subseteq \{Mh > \tfrac\delta2\} \cup \{\abs h >
\tfrac\delta2\}$ tiene \omterm{def:b3:measure:measure}{medida}
$\leq \frac{6\varepsilon}\delta
+ \frac{2\varepsilon}\delta$ (pregunta 12; Markov). $\varepsilon$
arbitrario: $\lambda(\Omega_\delta) = 0$; y uniendo sobre
$\delta = \frac1k$: $A_rf \to f$ en casi todo punto.

\textbf{16.} (a) Para cada $q \in \Q$, la pregunta 15 aplicada a
$\abs{f - q}$ da $A_r\abs{f - q}(x) \to \abs{f(x) - q}$ en casi todo
punto; sobre la intersección de estos conjuntos de
\omterm{def:b3:measure:measure}{medida} total, elíjase $q$ con
$\abs{f(x) - q} < \eta$: $\limsup_rA_r\abs{f - f(x)}(x) \leq 2\eta$ para
todo $\eta$: casi todo $x$ es un punto de Lebesgue. (b) En un punto de
Lebesgue,
\[
\Bigl|\frac{F(x + h) - F(x)}h - f(x)\Bigr|
= \Bigl|\frac1h\int_x^{x+h}\bigl(f - f(x)\bigr)\Bigr|
\leq 2\,A_{\abs h}\abs{f - f(x)}(x) \to 0 :
\]
$F' = f$ en casi todo punto --- las primitivas de funciones $L^1$ sí se
derivan de vuelta; la escalera (el~\cref{pb:b3:measure:1}) es
contraejemplo únicamente del sentido \emph{recíproco}.

\textbf{17.} Aplíquese la pregunta 15 a $\mathbf 1_{A\cap[-n,n]}$ y
hágase crecer $n$: para casi todo $x \in A$, la
\omterm{ex:b3:lebesgue:gamma}{densidad}
$\frac{\lambda(A\cap\intcc{x-r}{x+r})}{2r} \to 1$. Steinhaus: alrededor
de un punto de \omterm{ex:b3:lebesgue:gamma}{densidad}, tómese $r$ con
\omterm{ex:b3:lebesgue:gamma}{densidad} $> \frac34$; para
$\abs t < \frac r2$, $A$ y $A + t$ llenan cada uno más de $\frac32r$ de
un intervalo de longitud $\leq \frac52r$ y, por tanto, se cortan:
$\intoo{-r/2}{r/2} \subseteq A - A$.

\textbf{18.} Sea $G(x) = x^{\alpha+1-p}F(x)^p$: $G$ se anula en $0$ ($F$
se anula cerca de $0$) y en $\infty$ ($F$ acotada,
$\alpha + 1 - p < 0$), de modo que $\int_0^\infty G' = 0$ con
\[
G' = (\alpha + 1 - p)\,x^{\alpha-p}F^p +
p\,x^{\alpha+1-p}F^{p-1}f
= (\alpha + 1 - p)\Bigl(\frac Fx\Bigr)^px^\alpha +
p\Bigl(\frac Fx\Bigr)^{p-1}f\,x^\alpha .
\]
Por tanto $\int(F/x)^px^\alpha = \frac{p}{p-1-\alpha}
\int(F/x)^{p-1}f\,x^\alpha$; Hölder para la
\omterm{def:b3:measure:measure}{medida} $x^\alpha\dd x$ (exponentes
$\frac p{p-1}$ y $p$) remata como en la pregunta 3. Caso límite
$\alpha = p - 1$: con $f(t) =
\frac1t\mathbf 1_{\intcc1A}$, el miembro derecho es $\ln A$ mientras que
el izquierdo contiene $\int_1^A\frac{(\ln x)^p}x\dd x =
\frac{(\ln A)^{p+1}}{p+1}$: ninguna constante sobrevive a $A \to
\infty$.

\textbf{19.} Tonelli sobre $\{0 < t < x\}$:
\[
\langle Hf, g\rangle =
\int_0^\infty\frac{g(x)}x\int_0^xf(t)\,\dd t\,\dd x
= \int_0^\infty f(t)\int_t^\infty\frac{g(x)}x\,\dd x\,\dd t
= \langle f, H^*g\rangle .
\]
Dualidad: $\norm{H^*f}_p = \sup\{\langle f, Hg\rangle : g
\geq 0, \norm g_q \leq 1\} \leq \norm f_p\cdot
\sup\norm{Hg}_q \leq \frac q{q-1}\norm f_p = p\,\norm f_p$, invocándose
Hardy en $L^q$, cuya constante $\frac q{q-1}$ vale $p$.

\textbf{20.} Sepárese a lo largo de la diagonal (nula). En $\{y \leq
x\}$:
\[
\iint_{y\leq x}\frac{f(x)g(y)}{x}\,\dd y\,\dd x
= \int f(x)\,(Hg)(x)\,\dd x \leq \norm f_p\,\norm{Hg}_q
\leq p\,\norm f_p\norm g_q,
\]
puesto que Hardy en $L^q$ lleva la constante $\frac q{q-1} =
p$. Simétricamente, $\iint_{y>x} = \int g\,(Hf) \leq
\norm g_q\norm{Hf}_p \leq q\,\norm f_p\norm g_q$ ($\frac
p{p-1} = q$). En total: $(p + q)\,\norm f_p\norm g_q$.

\textbf{21.} Para $a_n = n^{-1/p}$, $n \leq N$: el miembro derecho es
$\bigl(\frac p{p-1}\bigr)^p\sum_{n\leq N}\frac1n =
\bigl(\frac p{p-1}\bigr)^p\ln N + O(1)$. En el izquierdo, para $n
\leq N$: $a_1 + \dots + a_n \geq \int_1^{n+1}t^{-1/p}\dd t =
\frac{p}{p-1}\bigl((n+1)^{1-1/p} - 1\bigr)$, de modo que la $n$-ésima
media de Cesàro es $\geq \frac p{p-1}n^{-1/p}(1 -
o(1))$ uniformemente para $n$ en cualquier rango $n \geq n_0(\eta)$;
elevando a $p$ y sumando, el miembro izquierdo es $\geq
\bigl(\frac p{p-1}\bigr)^p(1 - \eta)\ln N + O_\eta(1)$. Dividiendo y
haciendo $N \to \infty$, y después $\eta \to 0$: ninguna constante menor
que $\bigl(\frac p{p-1}\bigr)^p$ puede servir.

\textbf{22.} $H$ acotado en $L^p$: los promedios acumulados no inflan
las normas $p$ (constante $\frac p{p-1}$); Cesàro en $\ell^p$: lo mismo,
discretizado; $M$ acotado en $L^p$: incluso el mejor promedio local
queda bajo control (constante $O(\frac1{p-1})$). Los tres fallan en
$p = 1$, y por una sola razón: promediar una unidad de masa concentrada
produce una cola $\frac1x$ (preguntas 10 y 13), y $\frac1x$ pertenece
cerca del infinito a todo $L^p$ salvo a $L^1$. El suavizado esparce la
masa exactamente hasta la frontera armónica de la
\omterm{def:b3:lebesgue:l1}{integrabilidad}.

\textbf{23.} Póngase $b_n = a_n^{1/p}$, de modo que $\sum b_n^p = \sum
a_n$. La pregunta 8 da
\[
\sum_{n\geq1}\Bigl(\frac{b_1 + \dots + b_n}n\Bigr)^{p}
\leq \Bigl(\frac p{p-1}\Bigr)^{p}\sum_{n\geq1}a_n,
\]
y la desigualdad entre las medias aritmética y geométrica acota cada
sumando por abajo:
\[
\Bigl(\frac{b_1 + \dots + b_n}n\Bigr)^{p}
\geq \bigl(b_1\cdots b_n\bigr)^{p/n}
= \bigl(a_1\cdots a_n\bigr)^{1/n}.
\]
Por tanto, $\sum(a_1\cdots a_n)^{1/n} \leq \bigl(\frac
p{p-1}\bigr)^p\sum a_n$ para \emph{todo} $p > 1$. Con $m =
p - 1$,
\[
\Bigl(\frac p{p-1}\Bigr)^{p} =
\Bigl(1 + \frac1m\Bigr)^{m+1},
\]
que decrece hacia $\eu$ cuando $m \to \infty$ (la clásica sucesión
monótona superior de $\eu$). Tomando el ínfimo sobre $p$ se obtiene la
desigualdad de Carleman con constante $\eu$.

\textbf{24.} Para $a_n = \frac1n$, $(a_1\cdots a_n)^{1/n} =
(n!)^{-1/n}$. El encaje del~\cref{pb:b3:product:1} da
$\sqrt{2\pi n}\,(n/\eu)^n \leq n! \leq \sqrt{2\pi
n}\,(n/\eu)^n\eu^{1/(12n)}$, luego
\[
(n!)^{1/n} = \frac n\eu\,(2\pi n)^{1/(2n)}
\eu^{O(1/n^2)}
= \frac n\eu\Bigl(1 +
O\Bigl(\frac{\ln n}{n}\Bigr)\Bigr),
\]
puesto que $(2\pi n)^{1/(2n)} = \exp\bigl(\frac{\ln(2\pi
n)}{2n}\bigr)$. Invirtiendo, $(n!)^{-1/n} = \frac\eu n(1 +
O(\frac{\ln n}n))$, y sumando sobre $n \leq N$:
\[
\sum_{n\leq N}(n!)^{-1/n}
= \eu\ln N + O(1),
\qquad
\eu\sum_{n\leq N}\frac1n = \eu\ln N + O(1)
\]
(la serie de errores $\sum\frac{\ln n}{n^2}$ converge). Una desigualdad
de Carleman con constante $c$ obligaría a $\eu\ln N
+ O(1) \leq c\,(\ln N + O(1))$ y, por tanto, a $c \geq \eu$ tras dividir
por $\ln N$. Las sucesiones optimizadoras se alinean: la constante de
Hardy se aproxima con $n^{-1/p}$ (pregunta 21) y la de Carleman con
$n^{-1}$ --- en cada caso, la sucesión de tipo armónico que queda justo
fuera del espacio que se promedia.

\textbf{25.} Para $f = \mathbf 1_{\intcc01}$: $Hf(x) = 1$ en $\intoc01$
y $Hf(x) = \frac1x$ para $x > 1$, luego
$\norm{Hf}_2^2 = 1 + \int_1^\infty x^{-2}\dd x = 2$ y el cociente es
$\sqrt2 \approx 1.414$, aproximadamente el $71\%$ de la cota óptima.
Para $f_A$ ($p = 2$): $F(x) = 2(\sqrt x - 1)$ en $\intcc1A$, de modo que
en ese rango $Hf_A = 2x^{-1/2} - 2x^{-1}$ y, para $x > A$,
$Hf_A(x) = 2(\sqrt A - 1)/x$. Elevando al cuadrado e integrando,
\begin{align*}
\int_1^A\bigl(2x^{-1/2} - 2x^{-1}\bigr)^2\dd x
&= 4\ln A - 12 + \frac{16}{\sqrt A} - \frac4A,\\
\int_A^{\infty}\frac{4(\sqrt A - 1)^2}{x^2}\,\dd x
&= 4 - \frac{8}{\sqrt A} + \frac4A,
\end{align*}
de donde $\norm{Hf_A}_2^2 = 4\ln A - 8 + 8A^{-1/2}$; dividiendo por
$\norm{f_A}_2^2 = \ln A$ se obtiene la identidad enunciada. En $A
= \eu^{10}$: $4 - \frac{8(1 - \eu^{-5})}{10} = 3.2054$, luego el
cociente es $\sqrt{3.2054} \approx 1.790 < 2$. El defecto
$4 - \norm{Hf_A}_2^2/\norm{f_A}_2^2 \sim 8/\ln A$ decae solo
logarítmicamente: para alcanzar el cociente $1.99$ haría falta
$\ln A \approx 200$, es decir, $A \approx 10^{87}$. Las constantes
óptimas son teoremas, no experimentos.
\end{solution}
